%=====================================================================
\chapter{MLOps: Deployment, Monitoring and Reliability}\label{ch:mlops}
%=====================================================================
\locator{5}{Part V --- Data Science in Systems Context}

\begin{outcomes}
By the end of this chapter you will be able to:
\begin{tight}
  \item \textbf{Describe} the path from a trained model to a monitored production service.
  \item \textbf{Choose} between batch, online and edge serving for a stated requirement.
  \item \textbf{Specify} what to version so that a prediction can be reproduced months later.
  \item \textbf{Design} a monitoring plan covering data, prediction and outcome drift.
  \item \textbf{Plan} a safe rollout and rollback.
\end{tight}
\end{outcomes}

\begin{prereq}
\chref{cloudedge} (containers and deployment tiers), \chref{timeseries} (drift) and
\chref{evaluation} (the metrics you will monitor).
\end{prereq}

\begin{keyterms}
MLOps $\bullet$ model registry $\bullet$ feature store $\bullet$ training--serving
skew $\bullet$ batch / online / edge inference $\bullet$ CI/CD $\bullet$ shadow
deployment $\bullet$ canary release $\bullet$ A/B test $\bullet$ rollback $\bullet$
data drift $\bullet$ prediction drift $\bullet$ ground-truth lag $\bullet$
observability $\bullet$ model card
\end{keyterms}

\section{The Gap Between a Notebook and a System}

\begin{conceptbox}[What Changes at Deployment]
A notebook model runs once, on data you chose, with you watching. A deployed model
runs thousands of times a day, on data nobody inspected, with nobody watching, for
years --- while the world it was trained on changes. MLOps is the engineering
discipline that makes the second situation survivable.
\end{conceptbox}

\dsfig{%
\begin{tikzpicture}[font=\scriptsize,
  s/.style={rounded corners=3pt, draw=#1, fill=#1!8, text=#1!50!black,
            text width=2.35cm, align=center, font=\tiny\bfseries, minimum height=0.95cm},
  arr/.style={-{Stealth[length=2mm]}, primary!55, line width=1pt},
  note/.style={font=\tiny\itshape, text=gray!55!black, align=center, text width=2.6cm}]

\node[s=secondary] (a) at (0,0)     {1. TRAIN\\{\tiny versioned code + data}};
\node[s=goldhl]    (b) at (3.1,0)   {2. VALIDATE\\{\tiny gate on metrics}};
\node[s=tealhl]    (c) at (6.2,0)   {3. REGISTER\\{\tiny model registry}};
\node[s=primary]   (d) at (9.3,0)   {4. PACKAGE\\{\tiny container}};
\node[s=greenhl]   (e) at (12.4,0)  {5. ROLL OUT\\{\tiny shadow $\to$ canary}};
\node[s=purplehl]  (f) at (12.4,-2.3){6. MONITOR};
\node[s=accent]    (g) at (9.3,-2.3) {7. DETECT DRIFT};
\node[s=secondary] (h) at (6.2,-2.3) {8. RETRAIN};

\foreach \x/\y in {a/b,b/c,c/d,d/e}{\draw[arr] (\x) -- (\y);}
\draw[arr] (e) -- (f);
\draw[arr] (f) -- (g);
\draw[arr] (g) -- (h);
\draw[arr] (h.west) .. controls +(-2.0,0) and +(0,-1.4) .. (a.south);

\node[note] at (3.1,-1.15) {block promotion if it fails, don't just warn};
\node[note] at (12.4,-1.15) {never a big-bang release};
\node[note] at (3.1,-2.3) {the loop never terminates};

\node[anchor=north west, align=left, text width=13.8cm, font=\scriptsize, text=accent!75!black]
     at (-1.5,-3.4)
     {\textbf{Step 2 is the one teams omit.} An automated gate that refuses to promote a
      model scoring below the current production model --- on a fixed benchmark set,
      including fairness slices --- prevents the most common production incident: a
      retrained model silently worse than the one it replaced.};
\end{tikzpicture}}%
{The MLOps lifecycle. It is a closed loop, and the arrow from retraining back to training
is the one that distinguishes a maintained system from an abandoned
one.}{mlops-lifecycle}

\section{Serving Patterns}

\rowcolors{2}{white}{lightbg}
\begin{longtable}{@{}p{2.5cm}p{3.7cm}p{3.5cm}p{4.4cm}@{}}
\toprule
\rowcolor{primary!15}
& \textbf{Batch} & \textbf{Online (API)} & \textbf{Edge} \\
\midrule
\endhead
When it runs & On a schedule & On request & On the device, continuously \\
Latency & Hours & 10--500 ms & $<10$ ms \\
Complexity & Lowest & Moderate & Highest --- updates are hard \\
Good for & Weekly risk scores, nightly maintenance lists & Fraud checks, live recommendations & Safety shutdowns, offline operation \\
Watch out for & Staleness --- the score is as old as the last run & Latency budget, autoscaling, cold starts & Model size, power, over-the-air updates \\
\bottomrule
\end{longtable}

\begin{conceptbox}[Choose the Simplest That Meets the Requirement]
Most student projects reach for an API when a nightly batch job would serve the
decision perfectly well. If the intervention happens on Monday mornings, a Sunday
night batch is simpler, cheaper, easier to monitor and easier to reproduce.
Real-time serving is a cost you should incur only when the decision genuinely
cannot wait.
\end{conceptbox}

\section{Training--Serving Skew}

\begin{alertbox}[The Most Common Production Failure]
The model was trained on features computed by a pandas script, and serves on
features computed by a Java service written by a different team. The two
implementations disagree slightly --- a different null handling, a different
timezone, a rounding difference --- and accuracy in production is far below what
testing promised. Nothing is broken and no error is raised.

\textbf{The fix:} compute features in \emph{one} place. A \term{feature store}
serves the same computed features to both training and inference, which eliminates
the class of bug entirely rather than testing for it.
\end{alertbox}

\section{What to Version}

\begin{definitionbox}[Reproducibility Requires Five Things]
To reproduce a prediction made eight months ago you need: the \textbf{code}, the
\textbf{data} (or a snapshot reference), the \textbf{environment} (container image),
the \textbf{hyperparameters}, and the \textbf{model artefact} itself --- each tied
together by a single run identifier. Missing any one makes the result
unreproducible, and an unreproducible result cannot be defended when challenged.
\end{definitionbox}

\dsfig{%
\begin{tikzpicture}[font=\scriptsize,
  v/.style={rounded corners=3pt, draw=#1, fill=#1!10, text=#1!50!black,
            text width=2.3cm, align=center, font=\tiny\bfseries, minimum height=1.05cm}]

\node[v=secondary] (c) at (0,0)   {\faIcon{code}\\CODE\\{\tiny git commit}};
\node[v=goldhl]    (d) at (2.75,0){\faIcon{database}\\DATA\\{\tiny snapshot hash}};
\node[v=tealhl]    (e) at (5.5,0) {\faIcon{cube}\\ENVIRONMENT\\{\tiny image digest}};
\node[v=purplehl]  (p) at (8.25,0){\faIcon{sliders-h}\\PARAMS\\{\tiny config file}};
\node[v=greenhl]   (m) at (11.0,0){\faIcon{box}\\MODEL\\{\tiny artefact hash}};

\node[anchor=north, rounded corners=4pt, draw=primary, fill=primary!8, text=primary,
      text width=12.6cm, align=center, font=\scriptsize\bfseries, minimum height=0.62cm]
      (run) at (5.5,-1.5) {RUN ID 2026-08-01-a41f  ---  one identifier binds all five};
\foreach \n in {c,d,e,p,m}{
  \draw[-{Stealth[length=1.6mm]}, primary!45] (\n) -- (run);}

\node[anchor=north west, align=left, text width=13.4cm, font=\scriptsize, text=gray!55!black]
     at (-1.3,-2.5)
     {\textbf{Plus, for every individual prediction served:} the run ID of the model that
      produced it, the input features as actually received, the output, and the timestamp.
      Without this you cannot answer ``why was this student flagged in March?'' --- a
      question that will be asked, and which \chref{ethics} may require you to answer.};
\end{tikzpicture}}%
{The five things that must be versioned together. A model registry exists to keep these
bound to one identifier, so that any past prediction can be explained and
reproduced.}{versioning}

\section{Monitoring}

\dsfig{%
\begin{tikzpicture}[font=\scriptsize,
  hd/.style={rounded corners=3pt, fill=#1, text=white, font=\scriptsize\bfseries,
             text width=6.7cm, align=center, minimum height=0.55cm},
  b/.style={rounded corners=3pt, draw=#1, fill=#1!7, text width=6.7cm, align=left,
            font=\scriptsize, inner sep=6pt, minimum height=2.1cm}]

\node[hd=secondary] at (0,0)    {LAYER 1 --- INPUTS};
\node[b=secondary]  at (0,-1.42){%
  \textbf{Available:} instantly\\[3pt]
  Feature distributions vs.\ training; null rates; unseen categories; row volume.\\[3pt]
  \emph{Detects data drift and broken pipelines --- your earliest warning, and the only
  one that needs no labels.}};

\node[hd=goldhl] at (7.3,0)    {LAYER 2 --- PREDICTIONS};
\node[b=goldhl]  at (7.3,-1.42){%
  \textbf{Available:} instantly\\[3pt]
  Distribution of scores; alert rate; share above threshold; per-segment rates.\\[3pt]
  \emph{A model that suddenly flags 12\% instead of 3\% is telling you something before
  any outcome is known.}};

\node[hd=purplehl] at (0,-3.2)    {LAYER 3 --- OUTCOMES};
\node[b=purplehl]  at (0,-4.62){%
  \textbf{Available:} after a delay\\[3pt]
  True precision and recall, once ground truth arrives.\\[3pt]
  \emph{The number you actually care about, and the one you get last --- weeks for
  student outcomes, months for equipment failure.}};

\node[hd=greenhl] at (7.3,-3.2)    {LAYER 4 --- BUSINESS};
\node[b=greenhl]  at (7.3,-4.62){%
  \textbf{Available:} slowest\\[3pt]
  Stoppages prevented, students retained, breaches caught, cost saved.\\[3pt]
  \emph{The only layer the funder cares about. Report it even though it lags.}};

\node[anchor=north west, align=left, text width=13.4cm, font=\scriptsize, text=accent!75!black]
     at (-3.35,-6.5)
     {\textbf{Ground-truth lag is why layers 1 and 2 exist.} If you only monitor accuracy,
      you learn a model broke months after it broke. Input and prediction monitoring give
      you a signal on day one --- imperfect, but immediate. Every deployed model in this
      book should have all four layers, in that order of alerting priority.};
\end{tikzpicture}}%
{Four monitoring layers, ordered by how quickly each becomes available. Alert on layers 1
and 2; report layers 3 and 4.}{monitoring-layers}

\section{Safe Rollout}

\begin{definitionbox}[Shadow, Canary, and Rollback]
\term{Shadow}: the new model runs on live traffic and its predictions are logged
but not acted on --- zero risk, real data. \term{Canary}: it serves a small
percentage of real traffic, with automatic rollback if metrics degrade.
\term{Rollback}: reverting to the previous version, which must be a single command
and must be rehearsed.
\end{definitionbox}

\begin{worked}[Rolling Out a New Model Safely]
The dropout model is retrained and scores better offline. How does it reach
production?

\smallskip
\textbf{Week 1 --- shadow.} Both models score every student; only the old model's
output drives outreach. Compare on live data: do they agree? Where they disagree,
which is right? This catches training--serving skew, which offline testing cannot.

\textbf{Week 2 --- canary at 10\%.} A random 10\% of students are scored by the new
model. Automatic rollback triggers if the alert rate moves more than 50\% from
baseline, or if any monitored fairness slice degrades.

\textbf{Week 3 --- 50\%, then 100\%.} Advance only if layer-1 and layer-2 monitors
are stable.

\textbf{Week 8 --- outcome evaluation.} The first real precision figures arrive as
term results come in. Only now do you actually know whether the new model was
better.

\smallskip
\textbf{Keep the old model deployable for a full outcome cycle.} Offline metrics
promised an improvement; only week 8 confirms it, and until then rollback must
remain one command away.
\end{worked}

%---------------------------------------------------------------------
\begin{fourlenses}{Operating Models in Production}
\lensLA{\emph{Serving:} weekly batch --- outreach happens on a schedule, so an API is
unnecessary complexity. \emph{Ground-truth lag:} an entire term, which is why input
and prediction monitoring carry the alerting load. \emph{Special requirement:}
every prediction must be logged with its features and model version, because a
student may ask why they were contacted and the institution must be able to answer
(\chref{ethics}). \emph{Retraining:} annually, per cohort --- retraining mid-term
changes the rules on students partway through.}

\lensPM{\emph{Serving:} batch, weekly. \emph{Special difficulty:} the prediction
changes the outcome --- a team told it will miss the sprint often does not, which
makes the model look wrong precisely when it worked. Log the prediction \emph{and}
whether it was shown, so this can be analysed rather than argued about.
\emph{Retraining:} whenever team composition changes materially, not on a calendar.}

\lensIOT{\emph{Serving:} edge, which makes deployment the hard part --- over-the-air
updates across a fleet on unreliable links, with staged rollout by device group and
automatic rollback if a device fails to check in. \emph{Monitoring:} you must watch
the \emph{devices} as well as the model, and distinguish model drift from sensor
drift (\chref{iot}). \emph{Retraining:} monthly in the cloud on accumulated
features, pushed down after canary on a device subset.}

\lensSEC{\emph{Serving:} online and streaming --- detection latency is a control.
\emph{Retraining:} frequent, because the adversary adapts --- but each retrain is an
attack surface (\chref{cybersecurity}), so the training set must be screened and
promotion gated by a human. \emph{Monitoring:} alerts per analyst per hour is a
first-class operational metric, and a sharp \emph{drop} in alert volume is as
alarming as a spike, since it may mean a log source has silently died.}
\end{fourlenses}

%---------------------------------------------------------------------
\begin{lab}[A Monitored Service]
\begin{lstlisting}[language=Python]
import json, hashlib, numpy as np, pandas as pd
from datetime import datetime
from scipy.stats import ks_2samp

MODEL_VERSION = "2026.08.1"
TRAINING_REF  = pd.read_parquet("reference_features.parquet")   # training snapshot

# --- LAYER 1: input drift, available immediately -----------------------
def input_drift(live: pd.DataFrame, ref=TRAINING_REF, alpha=0.01):
    """KS test per feature. Note: with large n everything is 'significant'
       (Chapter 9), so gate on the STATISTIC, not the p-value."""
    report = {}
    for col in ref.columns:
        stat, _ = ks_2samp(ref[col].dropna(), live[col].dropna())
        report[col] = {"ks": round(float(stat), 4), "drifted": stat > 0.15}
    unseen = set(live.get("programme", [])) - set(ref.get("programme", []))
    report["_unseen_categories"] = list(unseen)
    report["_null_rates"] = live.isna().mean().round(3).to_dict()
    return report

# --- LAYER 2: prediction drift -----------------------------------------
def prediction_drift(scores, baseline_rate, threshold, tol=0.5):
    rate = float((scores > threshold).mean())
    return {"alert_rate": round(rate, 4),
            "baseline": baseline_rate,
            "drifted": abs(rate - baseline_rate) > tol * baseline_rate}

# --- every prediction is logged so it can be explained later ----------
def serve(features: pd.DataFrame):
    scores = model.predict_proba(features)[:, 1]
    for i, row in features.iterrows():
        log_prediction({
            "run_id":        MODEL_VERSION,
            "entity_id":     row["student_id"],
            "features":      row.drop("student_id").to_dict(),   # as RECEIVED
            "feature_hash":  hashlib.sha256(
                                 json.dumps(row.to_dict(), sort_keys=True,
                                            default=str).encode()).hexdigest()[:12],
            "score":         float(scores[i]),
            "ts":            datetime.utcnow().isoformat(),
        })
    return scores

# --- the promotion gate: refuse to ship a worse model -----------------
def promotion_gate(new_metrics, prod_metrics, slices):
    if new_metrics["pr_auc"] < prod_metrics["pr_auc"]:
        raise RuntimeError("BLOCKED: new model is worse on the benchmark set")
    for s in slices:                              # fairness slices, Chapter 28
        if new_metrics[s]["recall"] < prod_metrics[s]["recall"] - 0.05:
            raise RuntimeError(f"BLOCKED: recall regressed on slice '{s}'")
    return True
\end{lstlisting}
\end{lab}

%---------------------------------------------------------------------
\begin{checkpoint}
\begin{tightnum}
  \item Define training--serving skew and give a concrete example. What eliminates
        the whole class of bug?
  \item Your outcome labels arrive three months late. Which monitoring layers give
        you a same-day signal, and what can each actually detect?
  \item A retrained model scores better offline. Why is that insufficient reason to
        replace the production model immediately?
\end{tightnum}
\end{checkpoint}

%---------------------------------------------------------------------
\begin{chaptersummary}
\begin{tight}
  \item MLOps closes the loop: train, validate, register, package, roll out, monitor,
        detect drift, retrain.
  \item An automated promotion gate that blocks a worse model prevents the most
        common production incident.
  \item Choose the simplest serving pattern that meets the latency the decision
        genuinely needs.
  \item Training--serving skew is eliminated by computing features in one place, not
        by testing harder.
  \item Version code, data, environment, parameters and artefact under one run ID,
        and log every prediction with the version that made it.
  \item Monitor inputs and predictions for an immediate signal; outcomes and business
        impact arrive later but are what matter.
  \item Roll out via shadow then canary, keep rollback to one command, and keep the
        old model deployable for a full outcome cycle.
\end{tight}
\end{chaptersummary}

\begin{reviewq}
  \item \textbf{(MCQ)} Running a new model on live traffic without acting on its
        output is called: (a)~canary (b)~shadow (c)~rollback (d)~A/B testing.
  \item \textbf{(MCQ)} Which monitoring layer requires no ground truth?
        (a)~outcome (b)~business impact (c)~input distribution (d)~precision.
  \item \textbf{(Short)} List the five things that must be versioned to reproduce a
        prediction, and say what goes wrong if the environment is omitted.
  \item \textbf{(Short)} Explain ground-truth lag and its consequence for how you
        design alerting.
  \item \textbf{(Short)} Why can a sudden \emph{drop} in alert volume be as serious as
        a spike?
  \item \textbf{(Applied)} Write the monitoring plan for the predictive maintenance
        model of \chref{iot}: state what you measure at each of the four layers, the
        alert threshold for each, and who is paged.
  \item \textbf{(Applied)} Your team wants to auto-retrain and auto-deploy nightly.
        Give three arguments in favour, three risks, and the specific gates you would
        require before agreeing --- noting where the security domain differs.
\end{reviewq}
